% 英文 ABSTRACT
\thispagestyle{empty}
\begin{center}
  {\sffamily\zihao{-2}\bfseries ABSTRACT}
\end{center}

\vspace{1em}

With the rapid progress of large language models and tool-use techniques, AI agents are evolving from single-turn assistants into autonomous execution units that understand goals, decompose tasks, invoke external tools and continuously report progress. When multiple agents collaborate, a platform must simultaneously address task orchestration, inter-process communication, state consistency, execution isolation, real-time observability and fault recovery. Existing multi-agent frameworks are primarily built on Python platforms with role-driven dialogue flows, leaving room for improvement in local-first, high-performance and auditable distributed communication infrastructure. In particular, the Agent2Agent (A2A) protocol currently lacks an available SDK in the Rust ecosystem, which constrains the deployment of Rust-based local-first agent platforms.

This thesis designs and implements Fuxi, a Rust-based platform that addresses these gaps. Fuxi adopts a hierarchical Xuannv-Menke collaboration pattern that separates the user-facing orchestrator agent from worker agents. It provides the first from-scratch Rust implementation of A2A v1.0 compatible wire protocol, JSON-RPC server and client, and core data structures such as AgentCard, Task, Message and Artifact, with an additional InputRequired state for human-in-the-loop scenarios. A real-time event bus is built on top of Tokio broadcast and SQLite WAL so that live subscription and historical replay coexist on the same abstraction. Observability is exposed through Firehose, WebSocket, SSE and IM APIs, and execution is isolated through Git worktrees and a three-layer sandbox. The thesis further formalizes end-to-end latency decomposition, broadcast event latency lower bounds and a Worker selection function, providing a quantifiable basis for performance analysis.

Experimental evaluation on an Apple Silicon platform measures throughput, latency, parameter sensitivity and event bus stress along multiple dimensions. Fuxi sustains 654.66 tasks per second with 8 workers under 10 ms simulated jobs, and scales linearly to 1336.68 tasks per second with 16 workers at 83.5 percent scaling efficiency. The p50 task dispatch latency is 32.97 ms and the p99 cross-node event flow latency is as low as 0.15 ms. With 64 concurrent subscribers, the event bus sustains 10\,000 events per second with zero drops. These results demonstrate that combining event-driven semantics with append-only logging delivers a communication substrate that simultaneously achieves throughput, real-time responsiveness and traceability for local-first AI agent collaboration, providing a reusable foundation for future multi-agent system engineering.

\vspace{1.5em}

\noindent {\bfseries Key words}: AI Agent; Distributed Communication; Multi-Agent System; Event Bus; A2A Protocol
